\section{Related Work}
\label{sec:related-work}

Hebrew acronym interpretation has an established research history. HAADS addresses
abbreviation disambiguation in Hebrew and Aramaic Jewish Law documents
\citep{hacohenkerner2010haads}. For Modern Hebrew,
\citet{jacobs2020acronyms} develop an acronym dictionary from unannotated text and
enrich it with contextual information for disambiguation. Their approach also
addresses cases in which an expansion is absent from the acronym's document.
These studies establish that both collecting possible expansions and choosing
among them are substantive problems. Our study starts with an identified target
occurrence and an existing candidate inventory. It therefore evaluates expansion
and selection under that input contract, rather than end-to-end acronym detection
or dictionary construction. The quality and coverage of the supplied inventory
remain part of the conditions under which a selection score can be interpreted.

A related distinction concerns what an evaluation split measures.
\citet{chen2023gladis} introduce GLADIS, an English acronym benchmark covering
general, scientific, and biomedical domains. Its task splits keep training acronym
types out of validation and test. Their AcroBERT model jointly encodes a context
and a candidate expansion, using a triplet objective to distinguish matching
from nonmatching pairs. This provides a precedent for scoring candidate text and
transferring across task-training types. Our DictaBERT ranker uses pairwise binary
classification, so it is not a reproduction of AcroBERT's training procedure.

Data construction also affects this comparison. GLADIS builds its general and
biomedical evaluation sets by replacing annotated entity long forms with acronyms,
while repartitioning an existing acronym dataset for the scientific domain
\citep{chen2023gladis}. Replacement creates a labeled example without requiring a
naturally occurring acronym mention. However, the sentence was originally written
with the full expression. For our study, this distinction motivates retaining the
construction history of each example and distinguishing largely substituted
training material from natural development occurrences. It does not establish that
one source is necessarily easier, since the acronym and expansion distributions
also differ.

Generation introduces a separate evaluation problem: the answer may express the
right expansion without matching the reference string.
\citet{agrawal2022large} prompt an LLM to expand abbreviations in clinical snippets
without showing answer choices. They then map its output to the candidate with
the greatest contiguous character overlap. Their design separates the information
available to the generator from the inventory available to the evaluator. It also
shows why a reported generation score depends on the answer resolver. We do not
assume that this English clinical matching rule is appropriate for Hebrew, where
acceptable wording and the distinction between expanding a name and identifying
its referent require explicit treatment.

The same concern appears in word sense disambiguation.
\citet{meconi2025large} study selection among supplied definitions and separately
ask LLMs to generate definitions, explanations, and examples without a predefined
answer list. Their generation study uses expert judgments rather than treating
literal agreement with a dictionary definition as sufficient. Word definitions
and acronym expansions are different outputs, but the methodological lesson is
relevant: the two response formats require a stated basis for comparison. Our
selection metrics are implemented, while a final human judgment protocol for
free expansions remains to be completed.
